
\documentclass{article}

\usepackage{amsmath}
\usepackage{amssymb}
\usepackage{palatino}
\usepackage{fancyhdr}
\usepackage{fullpage}

% document size parameters

% document parameters
\setlength{\textheight}{9in}

\setlength{\parindent}{0.in}
\setlength{\parskip}{0.1in}

\pagestyle{fancyplain}
\setlength{\topmargin}{-0.4in}
\setlength{\headsep}{0.4in}
\lhead{\large{\textbf{MATH 232: Numerical Sol'n of Differential Eq's
II}}--Syllabus}
\rhead{\textit{Fall 2008}}


\begin{document}
\sloppy 
\textbf{Instructor:} 
Michael Sprague, Assistant Professor, School of Natural Sciences\\ 
Office: Science \& Engineering Rm.~332, 209-228-4179\\
Mobile: (209) 217-7240\\
E-mail: msprague@ucmerced.edu

\textbf{Lecture Time \& Location:} MW 2:00 pm -- 3:15 pm, COB 261

\textbf{Office Hours:} Monday \& Wednesday:  3:15--4:00 pm  (S\&E 332)\\
Alternatively, students are welcome to e-mail, call, or just stop by my
office.  I am rarely on campus on Tuesdays \& Thursdays, but can usually be
reached via email or cell phone.

\textbf{Course Webpage:} Resources and homework assignments are located at
\vspace{-0.15in}
\begin{center}
   \textbf{http://faculty.ucmerced.edu/msprague/math232\_f2008.html}
\end{center}

\textbf{Prerequisite:} Math 231 or Consent of Instructor 

\textbf{Course Units:} Four

\textbf{Catalog Description:} 
Fundamental methods presented in Math 231 are used as a base for discussing
modern methods for solving partial-differential equations.  Numerical
methods include variational, finite element, collocation,  spectral, and
FFT. Error estimates and implementation issues will be discussed.  A
significant amount of programming will be required.

\textbf{Course Goals:}
%
We will discuss three broad classes of numerical methods that are very
common in science and engineering today  for finding approximate solutions
to partial-differential equations (PDEs).  These methods are  ($i$)
finite-difference methods, ($ii$) spectral methods, and ($iii$)
finite-element methods.  These methods will be applied to standard elliptic,
parabolic, and hyperbolic PDEs.  At the end of this course, students will 

\vspace{-0.1in}
\begin{enumerate}

\item understand the fundamental theoretical underpinnings of the methods
discussed

\item understand the relative strengths and weaknesses of the methods in the
context of the three PDE classes

\item understand the importance of validation of numerical methods

\item understand the concepts of consistency, stability, and convergence

\item be able to effectively implement the numerical methods discussed

\end{enumerate}

\textbf{Learning Outcomes:}
At the end of this course, students should be able to 

\vspace{-0.1in}
\begin{itemize}

\item describe typical solution behavior for the three different PDE types
(supports goals 2, 5);

\item classify a PDE, and choose an appropriate and efficient numerical method
to find an approximate solution (supports goals 1--3, 5);

\item write a computer program than implements the numerical method
(supports goals 1, 3, 5); 

\item choose an appropriate discretization level for the desired accuracy
(supports goals 1, 4, 5);

\item demonstrate the validity of an approximate numerical solution
(supports goals 3, 5);

\item present effectively numerical results in written and graphical form
(supports goals 1--3).

\end{itemize}

\clearpage 
\textbf{Assessment \& Scoring:} 
%
The course grade is based on homework (85\%) and class participation (15\%).
This point distribution was chosen such that, for a student to earn an A, he
or she must participate in lectures by asking and answering questions.    If
a student has overall unsatisfactory class participation, I will email the
student asking for more involvement.  Please let me know if
you have any questions regarding your participation level.

Each homework assignment will be scored on a scale of 100 points.  Letter
grades will be given in the \textbf{approximate} framework: A: 90-100\%, B:
80--90\%, C: 70--80\%, D: 60--70\%.  Homework will be assigned approximately
every two weeks with a total of six or seven assignments.  Late homework
will not be accepted without prior permission.  

\textbf{Homework Requirements:} 
%
When preparing your homework, present your answers in a manner that
demonstrates your understanding of the concept at hand.   To this end, try
to understand the point behind me asking you the question or presenting you
with the task.  If you don't understand the point, please come talk to me
and we can discuss it further.

Homework presentation is expected to be of high quality in content,
technical writing, and presentation.   Please heed the following
when preparing homework:

\vspace{-0.1in}
\begin{itemize}

\item Homework should be computer generated (i.e., prepare your homework
with Word or LaTeX;  If you are not already familiar with the use of LaTeX,
I highly encourage you to use this course as an opportunity to learn it).

\item Present data graphically whenever possible. 

\item All graphs \& figures need to be numbered with appropriate captions
and properly annotated (legends, labeled axes, etc).  

\item Write up conclusions that you draw from any graphs and/or figures.

\item All programs should be included as appendices to the homework.  

\item Please try to layout your homework so that paper use is minimized. For
example, don't use a whole sheet of paper to print a three-line code.

\item You are encouraged to work in groups. However, \textbf{all work turned
in must be your own.} You must \textbf{list explicitly any outside sources
employed} (\textit{e.g.}\ websites, \textit{Mathematica}, books, etc.) for
each problem you solve.  This does not mean that you are allowed to copy a
solution should you find it posted elsewhere (see Academic integrity below).
All programs must be written by you; no ``copying and pasting.''

\end{itemize}

\textbf{Textbooks:} There are no required textbooks for this course.
However, reading assignments will be taken from the list of titles below,
which are held on a non-circulating basis in the library. 

\vspace{-0.1in}
\begin{itemize}

 \item Pinchover \& Rubinstein, \textit{An Introduction to Partial
Differential Equations}, Cambridge University Press, 2005

  \item Strikwerda, \textit{Finite Difference Schemes and Partial
Differential Equations}, SIAM, 2004

\item Canuto et al., \textit{Spectral Methods: Fundamentals in Single
Domains}, Springer, 2006

\item Elman, Silvester \& Wathen, \textit{Finite elements and fast iterative
solvers: with applications in incompressible fluid dynamics}, Oxford
University Press, 2005.  (The first chapter is available as a free online
sample; chapter 1 is all we may cover:
http://www.oup.co.uk/pdf/0-19-852868-X.pdf)

\end{itemize}

\clearpage
\textbf{Other Useful References:} 

\vspace{-0.1in}
\begin{itemize}

  \item Boyd, \textit{Chebyshev and Fourier Spectral Methods: Second Revised
Edition}, Dover, 2001

  \item Gottlieb \& Orszag, \textit{Numerical Analysis of Spectral Methods:
Theory and Applications}, SIAM, 1977

 \item Iserles, \textit{A First Course in the Numerical Analysis of
Differential Equations}, Cambridge University Press, 2000

 \item Oden \& Reddy, \textit{An Introduction to the Mathematical Theory of
Finite Elements} (\$16 version to be  
released by Dover on Aug 28)

\end{itemize}



\textbf{Computers, Software \& Programs:} You may use a calculator (graphing or
otherwise) or other computational tool (\textit{e.g.}, Mathematica, Maple,
Matlab, etc) to aid in your completion of homework assignments.  

Your homework will require a significant amount of programming and data
presentation.   While any programming language will do, I recommend either
Matlab or  free alternative to Matlab (Octave or python).    However, there
is at least one assignment where I will ask you to use the FFT feature in
Matlab.  When writing your programs, 

\vspace{-0.1in}
\begin{itemize}
\item use many subroutines;
\item write many comments.
\end{itemize}

\textbf{Dropping the Course:}  Please see the UC Merced \textit{General
Catalog} for more details.

\textbf{Special Accommodations:} If you qualify for accommodations
because of a disability, please submit a letter from Disability
Services to the instructor in a timely manner so that your needs may
be addressed.  Student Affairs determines accommodations based on
documented disabilities.
I will make every effort to accommodate all students who, because of
religious obligations, have conflicts with scheduled exams, assignments, or
required attendance.  Please speak with me during the first week of classes
regarding any potential academic adjustments or accommodations that may
arise due to religious beliefs during this term.

\textbf{Academic Integrity:}  Academic integrity is the foundation of an
academic community and without it none of the educational or research goals
of the university can be achieved.  All members of the university community
are responsible for its academic integrity.  Existing policies forbid
cheating on examinations, plagiarism and other forms of academic dishonesty.
The current policies for UC Merced are described in the  Academic Honesty
Policy.  Go to 
http://studentlife.ucmerced.edu/ and look under ``Student Judicial
Affairs.''


\end{document}
